
\subsubsection{6.1 Summary of Findings}

Three principal findings emerge from these experiments:

First, within the Pythia 1.4B-6.9B family, alignment-induced miscalibration is consistent with H2 (decision-boundary restructuring) rather than H1 (scale inflation). All nine Spearman rho values fall below the H1 threshold. This implies that post-hoc correction methods designed under an H1 assumption -- including temperature scaling and ATS variants -- may not address the actual mechanism. Temperature scaling adjusts confidence magnitude on a fixed answer distribution; H2 restructuring changes the distribution itself.

Second, DPO produced larger calibration degradation than PPO across all three model sizes. This observation reverses the pre-registered ordering prediction. A possible explanation is that DPO's token-level direct preference reshaping operates more aggressively on answer distributions than PPO's KL-constrained sequence-level optimization. However, this interpretation should be treated as exploratory given that the fallback checkpoints may not be training-equivalent. At 2.8B, the DPO-PPO difference falls within overlapping bootstrap confidence intervals and is not statistically distinguishable.

Third, H3 (framing susceptibility) is not supported in the softmax-ECE evaluation setting. Calibration degradation is domain-general within the Pythia family, not driven by adversarial framing.

\subsubsection{6.2 Implications for Post-Hoc Calibration Correction}

If the dominant miscalibration mechanism is boundary restructuring rather than scale inflation, then corrections that merely rescale confidence on the aligned model's preferred answers may be insufficient. Under 1.4B-PPO, 99.7\% of MMLU items have a different top-ranked option after alignment. Rescaling confidence on these new preferences does not address the underlying problem that the model has learned to prefer different -- and potentially incorrect -- answers.

Xie et al. (2024) reported that ATS reduces LLaMA-2-Chat ECE by 58-82\%. If alignment shifts decision boundaries in hidden-state space, ATS may succeed by learning per-input temperature adjustments that partially track boundary redistribution. This interpretation predicts that ATS effectiveness should correlate with the degree of H2 boundary shift. Testing this prediction was planned as H-M4 but was not executed.

\subsubsection{6.3 Scale Effects}

The 6.9B-DPO model's rho = 0.875 -- below but near the H1 threshold of 0.90 -- raises the possibility that the dominant mechanism may shift from H2 toward H1 at larger scales. If models above approximately 13B parameters exhibit rho >= 0.90 under DPO alignment, standard temperature scaling would be expected to become more effective. Testing this hypothesis requires applying the Spearman rho diagnostic to larger model families such as LLaMA-2 or Mistral.

\subsubsection{6.4 Limitations}

\textbf{Fallback checkpoints.} Public fallback checkpoints were used instead of the primary RLHFlow Pythia alignment checkpoints. These share the same base model and approximate training regime, but training equivalence is not guaranteed. The H2 mechanism finding (8/9 Spearman rho values below 0.85 with zero above 0.90) would require implausibly consistent checkpoint artifacts to overturn. The DPO >= PPO ordering, however, is more sensitive to checkpoint training differences and should be interpreted with caution.

\textbf{Single model family.} All results are restricted to Pythia 1.4B-6.9B. This family provides the cleanest controlled experiment but is a research model family not widely deployed in production. Whether H2 dominates in LLaMA-2, Mistral, or other families is an open empirical question.

\textbf{Scale range.} The scale range of 1.4B-6.9B is limited. No publicly available Pythia variants exceed 12B parameters, preventing resolution of the potential H1/H2 mechanism transition suggested by the 6.9B-DPO result.

\textbf{H-M4 not executed.} The planned experiment testing whether ATS corrects H2-type boundary shifts was not conducted. The claim about ATS correctability is therefore framed as motivated future work rather than an empirical finding.

\textbf{Shot-count asymmetry.} The H3 diagnostic compares 4-shot MMLU with 0-shot TruthfulQA. If few-shot prompting amplifies MMLU calibration degradation relative to 0-shot, this could contribute to the observed $\Delta ECE_MMLU$ > $\Delta ECE_TruthfulQA$ pattern independently of framing susceptibility.

\subsubsection{6.5 Broader Impact}

If RLHF-aligned models are miscalibrated via boundary restructuring rather than confidence inflation, standard reliability checks that probe whether a model's confidence on its predicted class is accurate may not capture the actual problem: that the predicted class itself has changed systematically. Alignment-induced answer restructuring may be more consequential than overconfidence on stable predictions, as it implies the model has learned different answer preferences across a wide range of inputs.
